\begin{textsample}{POS Dim 1 – vem\_ed\_19 – Score 84.00 – vem\_ed\_19/t095.txt}  \label{ex:f1_pos_043}
VEm Newsletter dos Projetos Colaborativos Internacionais NESTA EDIÇÃO ENTREVISTA com Rosi León, diretora de Intercâmbios Virtuais e Aprendizagem On-line na DePaul University (EUA) EVENTOS INTERNACIONAIS debatem gestão de projetos COIL (Collaborative Online International Learning) FATEC FRANCO DA ROCHA realiza PCIs na área de sustentabilidade Sustentabilidade \textbf{é} tema de PCIs Boas Práticas No primeiro semestre de 2023, a Fatec Franco da Rocha desenvolveu dois Projetos Colaborativos Internacionais (PCIs / Cesu) abordando a sustentabilidade.
A parceria \textbf{foi} com a professora Malgorzata Marchewka, da UEK (Universidade de Economia de Cracóvia, Polônia). \textbf{São} eles: Objetivos de Desenvolvimento \textbf{Sustentável} (ODS) da ONU Este PCI envolveu as \textbf{disciplinas} de Gestão de Projetos, ministrada pelo professor João Marcos Silva de Almeida, e Inglês V, ministrada pela professora Lúcia Maria dos Santos, ambas do curso de Gestão de Tecnologia da Informação.
Energias renováveis e negociações - PCI idealizado por Paulo Hélio Kanayama, diretor da Unidade, e desenvolvido pela professora Marcia Capelini, \textbf{que} leciona Energia e Ambiente no curso de Gestão de Energia e Eficiência Energética.
Neste projeto, o \textbf{apoio} \textbf{linguístico} ficou a cargo do professor Jorge Tenório Fernando.
Desde 2022, Jorge trabalha em parceria com Gosia, como a docente polonesa \textbf{é} mais conhecida.
O primeiro PCI entre a Fatec Franco da Rocha e a UEK \textbf{foi} notícia no site Cesu: https://cesu.cps.sp.gov.br/pci-entre-fatec-franco-da-rocha-e-uek-polonia-gera-live-para-angariar-doacoes-para-criancas-e-jovens-da-ucrania/ A \textbf{seguir}, confira os depoimentos de docentes e gestores envolvidos nos dois projetos \textbf{desenvolvidos} entre UEK e Fatec Franco da Rocha entre março e junho de 2023.
Suas \textbf{reflexões} sobre o \textbf{desenvolvimento} de \textbf{competências} por meio de PCIs voltados à discussão sobre sustentabilidade demonstram a \textbf{importância} da parceria internacional, do \textbf{apoio} da Direção e das Coordenações de Cursos e do trabalho em sinergia entre professores de disciplinas técnicas e de \textbf{idiomas}.
Małgorzata Marchewka professora da UEK (Universidade de Economia de Cracóvia, Polônia) Małgorzata Marchewka professora da UEK (Universidade de Economia de Cracóvia, Polônia) Neste semestre, realizamos dois projetos de Intercâmbio Virtual entre a Fatec Franco da Rocha e a UEK.
Com sucesso, pois os alunos não apenas adquiriram \textbf{conhecimentos} na área de gerenciamento de projetos, fornos solares e negociações.
Seus benefícios mais importantes incluem o aprimoramento das \textbf{habilidades} \textbf{linguísticas}, literacia digital e um aumento \textbf{geral} da autoconfiança.
Para muitos deles, essa \textbf{foi} a primeira experiência real de trabalho em equipes virtuais internacionais, em \textbf{que} conhecer a gramática inglesa não \textbf{é} suficiente para lidar com as \textbf{tarefas}.
O \textbf{que} se deseja \textbf{é} tolerância, paciência e abertura para os outros.
Eles deram o melhor de si e os resultados mostram isso!
Paulo Hélio Kanayama diretor da Fatec Franco da Rocha O \textbf{que} mais me chamou a atenção \textbf{foi} o reconhecimento da \textbf{importância} da sustentabilidade demonstrado pelos estudantes, \textbf{que} transmitem vontade de \textbf{mudanças} em relação ao paradigma de \textbf{desenvolvimento} vigente.
Esse \textbf{era} um dos objetivos do projeto: provocar engajamento da comunidade.
A segunda \textbf{coisa} \textbf{que} achei impressionante \textbf{foi} a desenvoltura dos brasileiros expressando-se em inglês.
Eles se esforçaram bastante, sem medo de \textbf{serem} julgados.
Esse \textbf{era} outro dos objetivos do PCI, \textbf{utilizar} o inglês como instrumento de comunicação.
João Marcos Silva de Almeida professor de Gestão de Projetos na Fatec Franco da Rocha Os alunos aplicaram os principais conceitos relativos às boas práticas em gestão de projetos: planejamento, escopo, gestão do tempo e da comunicação.
Entre os \textbf{aprendizados} estão o aprimoramento de \textbf{competências} como liderança, motivação, \textbf{flexibilidade} e trabalho em equipe.
A inclusão do tema ODS contribuiu para \textbf{que} os alunos refletissem sobre gestão de projetos e sustentabilidade.
Desejo \textbf{que} este PCI continue para \textbf{que} mais alunos sigam superando desafios e vivenciando na prática o projeto.
Para nós professores, também \textbf{é} uma excelente oportunidade de \textbf{interação} acadêmica internacional. coordenadora do curso de Gestão de Tecnologia da Informação da Fatec Franco da Rocha Queria manifestar minha satisfação com o PCI e enfatizar a \textbf{importância} dos aspectos relacionados à gestão de projetos e comunicação em inglês, também destaco o estudo dos ODS.
Essa abordagem \textbf{foi} importante, pois reforçou a pesquisa nos nossos projetos \textbf{interdisciplinares} do curso. professora de Energia e Ambiente na Fatec Franco da Rocha Sob a \textbf{perspectiva} \textbf{ambiental}, o PCI \textbf{foi} fundamental no \textbf{aprendizado} conjunto sobre a visão do “pensar globalmente” — uma vez \textbf{que} a busca de alternativas energéticas \textbf{sustentáveis}, diante dos desafios \textbf{que} \textbf{representam} as \textbf{mudanças} climáticas, supera as barreiras geográficas, \textbf{políticas} ou culturais — e do “agir localmente”.
Adaptar as soluções às \textbf{necessidades} \textbf{locais} \textbf{é} um grande desafio, \textbf{que} permeia os ODS.
Assim, ver esses jovens respondendo de \textbf{forma} tão criativa \textbf{foi} compensador.
Tenho certeza de \textbf{que} essa experiência irá agregar muito à formação \textbf{profissional} dos alunos, em um curso \textbf{que} já \textbf{é} diferenciado pela \textbf{integração} das \textbf{questões} \textbf{ambientais} como tema transversal nas \textbf{disciplinas}. professor de Inglês na Fatec Franco da Rocha A \textbf{turma} da Fatec, no \textbf{geral}, tinha \textbf{conhecimentos} elementares de inglês.
Por outro lado, dominava os aspectos técnicos do PCI.
No vídeo de apresentação final, os fatecanos explicaram o conceito dos fornos solares, suas vantagens e benefícios, ao passo \textbf{que} aos alunos da UEK, da \textbf{disciplina} de negociação, coube elencar um conjunto de argumentos sólidos para estimular sua aceitação e sua implementação como política pública.
As barreiras de comunicação \textbf{foram} superadas devido à \textbf{importância} do projeto e seu \textbf{impacto} positivo para os envolvidos, e, mais além, para a sociedade. @ Fale conosco Se você deseja desenvolver um Projeto Colaborativo Internacional (PCI) com alguma instituição de ensino estrangeira, preencha: https://forms.office.com/r/eH5ER2ZNKb Leitores Aos Osvaldo Succi Jr.
Coordenador PCIs Sustentabilidade \textbf{é} um tema \textbf{essencial} para a sobrevivência do planeta e para o ensino voltado aos Objetivos de Desenvolvimento \textbf{Sustentável} (ODS) da ONU, o \textbf{que} \textbf{é} uma tendência da educação.
No primeiro semestre de 2023, a Fatec Franco da Rocha realizou dois Projetos Colaborativos Internacionais (PCIs / Cesu): ODS e energias limpas, em parceria com a Universidade de Economia de Cracóvia (UEK, Polônia).
A matéria traz resumo dos dois projetos e trazem os depoimentos da professora polonesa Małgorzata Marchewka, do diretor da Fatec Franco da Rocha, Paulo Hélio Kanayama, da coordenadora do curso de Gestão de Tecnologia da Informação, Silvia Farani, e dos professores Jorge Tenório Fernando (Inglês), João Marcos Silva de Almeida (Gestão de Projetos) e Marcia Capelini (Energia e Ambiente).
Quando Direção, Coordenações de Curso, professores de \textbf{disciplinas} \textbf{profissionais} e de \textbf{idiomas} se unem pela Internacionalização em Casa, os resultados logo aparecem.
Confira também a entrevista com Rosi León, diretora de Intercâmbios Virtuais e Aprendizagem On-line na DePaul University (EUA).
Em 2023, a instituição completa 10 anos de Global Learning Experience (GLE), como lá \textbf{são} chamados os Intercâmbios Virtuais ou Collaborative Online International Learning (COIL).
Completando as \textbf{ações} da equipe de PCIs, veja a matéria sobre a apresentação de trabalhos acadêmicos sobre gestão de projetos COIL em dois eventos internacionais realizados em junho de 2023: Red LatAM COIL e COIL Connect.
Boa leitura!
Quem?
Quem \textbf{é} Rosi León, diretora de Intercâmbios Virtuais e Aprendizagem On-line na DePaul University (EUA) Rosi León \textbf{é} diretora de Intercâmbios Virtuais e Aprendizagem On-line na DePaul University (EUA).
Nascida na Bulgária, mudou-se para os Estados Unidos, onde concluiu a licenciatura em Educação de Adultos e Línguas Estrangeiras e a pós-graduação em Educação Bilíngue / Bicultural pela DePaul.
Além da língua materna e do espanhol, fala rápida e claramente o inglês de Chicago, cidade em \textbf{que} vive.
Seu sorriso reflete o entusiasmo \textbf{que} nutre pelos Intercâmbios Virtuais, chamados Global Learning Experience (GLE) na DePaul.
Rosi participa dessa experiência de aprendizagem global desde sua concepção, em 2013.
Entrevista Seção Conte como foi \textbf{implementar} GLE na DePaul e sua evolução ao \textbf{longo} da década.
Começamos com um \textbf{programa} de treinamento on-line para os docentes, com atividades predominantemente assíncronas e reuniões \textbf{síncronas} semanais.
Na primeira semana, os professores compartilham suas experiências; na segunda, fazem um brainstorming sobre ideias de projetos e, na \textbf{terceira}, abordamos o cenário intercultural, como planejar e prevenir \textbf{coisas} \textbf{que} podem dar errado ou, se estivermos em uma situação difícil, como gerenciar a experiência dos alunos e do professor.
Desde o \textbf{início}, fizemos parceria com o Center for Teaching and Learning, unidade da DePaul University \textbf{que} apoia professores com o design instrucional para cursos on-line, oferecidos desde muito antes da pandemia.
Essa equipe ajuda com aspectos tecnológicos e treinamento em tópicos de ensino-aprendizagem úteis para os professores.
O mais recente, por exemplo, \textbf{foi} sobre Inteligência Artificial.
Entre as diversas atividades do nosso Global Engagement Office está o treinamento de professores, o \textbf{que} \textbf{é} muito importante, especialmente para aqueles \textbf{que} estão \textbf{começando}.
Por muitos anos, nosso \textbf{foco} principal \textbf{foi} nos professores, pois sem eles não tocamos o \textbf{programa}.
Mais recentemente, \textbf{comecei} a me preocupar com a \textbf{forma} de chegar aos estudantes.
Muitos deles não \textbf{sabem} \textbf{que} vão participar de um Intercâmbio Virtual até o professor \textbf{contar}.
Então, estamos trabalhando em um \textbf{pequeno} módulo de preparação com informações básicas sobre o \textbf{que} \textbf{é} COIL / GLE.
Qual \textbf{foi} o pedido de parceria mais desafiador?
Frequentemente, os projetos mais desafiadores \textbf{são} aqueles com \textbf{disciplinas} muito \textbf{específicas}, como aconteceu com um professor de Matemática — seguimos buscando parceria.
Quando os docentes \textbf{são} novos na abordagem, há muito para \textbf{produzir} e planejar.
Depois \textbf{que} ganham experiência, percebem os benefícios de projetos \textbf{interdisciplinares}.
Seguimos \textbf{aprendendo} sobre o difícil \textbf{processo} de encontrar parceiros (matchmaking).
Há muitas \textbf{coisas} a considerar e \textbf{que} precisam \textbf{ser} resolvidas.
Uma delas \textbf{é} pedir informações \textbf{específicas}: breve descrição da \textbf{disciplina}, plano de ensino, tamanho típico da \textbf{turma} e qual o semestre ou período no qual \textbf{é} oferecida. \textbf{Que} conselho você daria para professores interessados em manter Intercâmbios Virtuais por um \textbf{longo} tempo? \textbf{Que} conselho você daria para professores interessados em manter Intercâmbios Virtuais por um \textbf{longo} tempo?
Tentamos fazer os professores enxergarem oportunidades \textbf{que} os projetos COIL oferecem para além da duração das colaborações.
Se \textbf{é} uma parceria nova, encorajamos a oferecer o mesmo projeto por vários semestres.
O planejamento da primeira edição leva mais tempo, mas da segunda oferta em diante, \textbf{são} \textbf{mudanças} \textbf{menores}.
Utilizar a mesma \textbf{estrutura} de projeto com outros grupos de estudantes, e com parceiros em diversos países, traz múltiplas \textbf{perspectivas}.
Outro aspecto \textbf{que} vejo crescer \textbf{é} a pesquisa: professores trabalhando juntos investigam a metodologia de Intercâmbios Virtuais em seu campo de \textbf{conhecimento} e em suas \textbf{turmas}.
E \textbf{publicam} artigos em periódicos acadêmicos sobre essa inovação.
Se têm parceiros em vários países, podem realizar pesquisas multidisciplinares em diversas regiões do mundo.
Há muita demanda para pesquisa sobre Intercâmbios Virtuais: sobre o \textbf{desenvolvimento} de \textbf{habilidades} dos alunos, sobre a aprendizagem sob as \textbf{perspectivas} comparadas das \textbf{turmas} e como isso muda de \textbf{acordo} com o país — os resultados \textbf{são} parecidos ou diferentes?
Há tantas oportunidades para \textbf{produzir} artigos, periódicos ou até mesmo capítulos de livros.
Qual sua visão para os Intercâmbios Virtuais nos próximos cinco anos?
Talvez por causa de meu \textbf{papel}, vejo \textbf{todos} os benefícios \textbf{que} os Intercâmbios Virtuais podem trazer.
Acho \textbf{que} COIL veio para ficar e vai crescer à medida \textbf{que} mais instituições se juntarem a essa comunidade.
A pandemia ajudou de certa \textbf{forma} as instituições \textbf{que} buscam oferecer experiências de internacionalização para seus estudantes. \textbf{Parte} de minha visão e esperança \textbf{é} \textbf{que} haja mais \textbf{apoio} estável para iniciativas de Intercâmbio Virtual nas instituições.
Muitas \textbf{começam} sem equipe, apenas com uma \textbf{pessoa} \textbf{que} vê COIL como uma grande oportunidade para seus professores, estudantes e para a instituição como um \textbf{todo}.
Depois de um tempo, conseguem mais recursos. \textbf{É} importante \textbf{reconhecer} o \textbf{papel} das \textbf{estruturas} de suporte nas instituições — tecnológico, de design de projetos, de parcerias, \textbf{linguístico}.
Acho \textbf{que} o \textbf{desenvolvimento} da tecnologia de Inteligência Artificial para auxiliar nas traduções simultâneas deve facilitar o trabalho \textbf{colaborativo}, especialmente o de \textbf{pequenos} grupos de estudantes.
Também preciso \textbf{mencionar} as redes de coordenadores de Intercâmbios Virtuais nas diferentes instituições. \textbf{É} importante \textbf{aprender} uns com os outros, porque há um \textbf{conhecimento} incrível nas diferentes regiões do mundo, e manter vivas essas redes. \textbf{É} uma comunidade maravilhosa, \textbf{que} pode se beneficiar de recursos e ferramentas \textbf{que} ajudem a \textbf{todas} as instituições.
Intercâmbios Virtuais em debate Matéria No fim do primeiro semestre de 2023, a equipe dos PCIs marcou presença on-line em dois eventos internacionais sobre gestão de projetos COIL (Collaborative Online International Learning): o 3º Congresso da Red LatAM COIL, realizado de 12 a 16 de junho, e o primeiro webinário de uma série organizada pela COIL Connect, em 14 de junho.
A fundação \textbf{é} liderada por Jon Rubin, pioneiro da abordagem COIL, \textbf{iniciada} na State University of New York em 2006.
Osvaldo Succi Junior, coordenador dos PCIs / Cesu, e Heather Ward, da University of North Carolina Chapel Hill (EUA), \textbf{compartilharam} suas experiências de liderança na \textbf{construção} de \textbf{programas} \textbf{bem-sucedidos} de Intercâmbios Virtuais em suas instituições e relataram a \textbf{expansão} desses projetos na comunidade COIL. “Leading COIL Initiatives in Different Environments \& Expanding the COIL Connect Community: A Discussion” \textbf{foi} o título da apresentação.
Após um breve \textbf{histórico} dos PCIs, a tropicalização dos Intercâmbios Virtuais \textbf{proposta} por Succi Junior e \textbf{iniciada} em 2013 na Fatec Americana, o coordenador dos PCIs trouxe um resumo dos principais resultados nesta década.
Entre 2014 e 2023, as Fatecs envolveram cerca de 7.500 alunos em PCIs.
Em média, desde 2020 \textbf{são} realizados 60 projetos por semestre.
A edição de VEm 17 trouxe os principais fatos e \textbf{números} desta primeira década de PCIs: https://cesu.cps.sp.gov.br/wp-content/uploads/2023/05/VEm-17.pdf Sustentabilidade das experiências O 3º Congresso da Red LatAM COIL \textbf{contou} com três apresentações da equipe de PCIs / Cesu.
No dia 12, o painel “La sostenibilidad de COIL antes y después de COVID: experiencias de bajo coste” tratou da reciclagem e estandardização de projetos de Intercâmbios Virtuais, da formação e capacitação de professores e redes de inovação docente \textbf{específicas} para COIL, bem como das redes de"embaixadores COIL"para \textbf{compartilhar} experiencias e divulgar os projetos interna e externamente a suas instituições.
Esse painel \textbf{foi} composto por Succi Junior com Eva Haug, consultora educacional na Amsterdam University of Applied Sciences (Holanda) e com coordenadoras de internacionalização de duas universidades públicas catalãs: Marina Vives (Universitat Rovira i Virgili, Tarragona) e Alicia Betts (Universidad de Girona).
Reflexões: gestão dos projetos Em 13 de junho, Ana Carolina Freschi, do \textbf{apoio} aos PCIs em inglês na equipe PCI / Cesu, apresentou o trabalho com seus colegas: Divinia Jithoo, Lindelwa Mkhize (Durban University of Technology, África do Sul), Daniel Otieno Okech (Kenyatta University, Quênia), Eva Haug (AUAS, Holanda) e Harshita Tripati (Shiv Nadar University, Índia).
Em pauta, a dinâmica de facilitação de Intercâmbios Virtuais nos diversos \textbf{contextos} educacionais desses países — sob os \textbf{pontos} de vista tecnológico, intercultural, acadêmico e \textbf{linguístico}.
Nesse mesmo dia, o painel “¡ SOS!
Intervención del Coordinador COIL” \textbf{foi} conduzido por Succi Junior, Gisselle Morales Veloquio (diretora de aprendizagem global do Tecnológico de Monterrey, México) e Regiane Moreira (professora responsável pelo \textbf{apoio} aos PCIs em espanhol nas Fatecs).
Os autores sugerem 10 passos para a intervenção do coordenador de Intercâmbios Virtuais, organizados nas seguintes etapas: institucionalizar; definir claramente o \textbf{papel} do coordenador COIL; \textbf{alinhar} \textbf{expectativas} da instituição, seus coordenadores, professores e parceiros; determinar critérios para busca de parceiros; oferecer formação continuada (para professores novos e experientes em Intercâmbios Virtuais); gerenciar o design dos projetos; realizar o acompanhamento dos projetos; \textbf{criar} índices de comparação e de informação; explicar e difundir constantemente os projetos COIL; decidir quando intervir (intervenção de emergência, apenas quando necessário). \textbf{continuação} Alguns sinais de \textbf{que} um projeto não vai bem: professores \textbf{que} não participam, não detalham o projeto; docentes \textbf{que} culpam o parceiro ou não buscam soluções / bons resultados para o projeto - considerado como atividade pontual sem continuidade e sem visão das relações interinstitucionais; professores \textbf{que} não escutam seus alunos, não \textbf{seguem} o planejado ou mudam os planos no meio do projeto sem consultar parceiros e não buscam um “justo meio-termo”; estudantes \textbf{que} não interagem, confusos ou em constante \textbf{conflito}, sem alguém para coordená-los, ou ainda insensíveis à cultura do outro país, vista sem respeito e com estereótipos.
Do \textbf{ponto} de vista \textbf{institucional}, portanto, \textbf{é} imprescindível \textbf{contar} com um coordenador \textbf{que} ajude a resolver os problemas \textbf{locais}.
Em resumo: para \textbf{expandir} os Intercâmbios Virtuais com qualidade, \textbf{é} preciso ter uma equipe \textbf{que} gerencie o design, a implantação e o \textbf{desenvolvimento} desses projetos.

% matched lemmas: acordo, alinhar, ambiental, apoio, aprender, aprendizado, ação, bem-sucedido, coisa, colaborativo, começar, compartilhar, competência, conflito, conhecimento, construção, contar, contexto, continuação, criar, desenvolvimento, desenvolvir, disciplina, específico, essencial, estrutura, expandir, expansão, expectativa, flexibilidade, foco, forma, geral, habilidade, histórico, idioma, impacto, implementar, importância, iniciar, institucional, integração, interação, interdisciplinar, início, linguístico, local, longo, mencionar, mudança, necessidade, número, papel, parte, pequeno, perspectiva, pessoa, político, ponto, processo, produzir, profissional, programa, propor, publicar, que, questão, reconhecer, reflexão, representar, saber, seguir, ser, sustentável, síncrono, tarefa, terceiro, todo, turma, utilizar
\end{textsample}
